%!TEX TS-program = xelatex
%!TEX encoding = UTF-8 Unicode
\documentclass[Linux]{xjtuBSThesis}

\renewcommand{\upcite}[1]{\cite{#1}}

%ai在写的时候注意要适当换行,不要都堆在同一行写


% ==================== 基本信息 ====================
\author{陈必纯}{Bin Chen}
\title{面向项目级数据理解的图谱化 Text-to-SQL 智能体设计与实现}{Design and Implementation of a Graph-Native Text-to-SQL Agent for Project-Level Data Understanding}
\advisor{XXX}{Prof. XXX}
\college{电信学部}
\department{自动化}
\class{自动化钱2201（信息）}
\studentid{2222112595}
\edate{2025/06/15}

\begin{document}

% ==================== 封面 ====================
\makecover

% ==================== 前置部分 ====================
\frontmatter

% ---------- 中文摘要 ----------
\begin{abstractcn}

自然语言数据分析正在从“生成一条 SQL”转向“让智能体理解并持续维护一个数据项目”。
在企业和科研场景中，关系数据库、表格文件、JSON 文档、项目说明、历史查询和业务规则往往共同构成分析上下文；
真实问题也不只依赖表名和列名，还涉及隐式连接关系、领域术语、指标口径、样例值和历史经验。
大语言模型具备较强的语言理解与代码生成能力，但在真实数据项目中仍面临上下文组织困难、schema/value grounding 不稳定、工具验证成本高以及知识难以跨问题复用等挑战。
针对这些问题，本文设计并实现 Pontis：一种面向项目数据理解的图谱化智能体工作空间。
Pontis 不以训练或微调专用 Text-to-SQL 模型为目标，而是在原始数据与大语言模型智能体之间构建可查询、可更新、可复用的项目知识层，
使智能体能够围绕数据资产逐步检索、验证和积累上下文，再完成自然语言查询与数据分析任务。

本文提出以知识图谱为核心的项目工作空间架构，将文件、数据库、表、列、外键、统计摘要、语义描述、消歧规则和经验知识等统一建模为图实体及关系。
系统通过自动化提取流程发布基础事实，并生成列统计、值分布、列重叠、外键验证、语义摘要和向量索引等派生知识；
同时通过智能体探索沉淀难以完全规则化的连接经验、字段消歧和任务先验。
在此基础上，Pontis 设计了面向数据分析的工具调用智能体，将图检索、元数据读取、SQL 查询、文本搜索和图写入等能力封装为工具，
并通过轮次限制、探索检查、SQL 实体检查、桥接表检查和 SQL 消歧检查等护栏机制，降低模型在复杂数据库问答中的幻觉、错误连接和过度探索风险。

【实验结果】

\end{abstractcn}

\keywordscn{数据智能体；Text-to-SQL；大语言模型；知识图谱}

% ---------- 英文摘要 ----------
\begin{abstracten}

As data projects in enterprise and research scenarios continue to grow, relational databases, tabular files, JSON documents, and project documentation often coexist in loosely organized directories. Large language models have shown strong capabilities in natural language understanding and code generation, but they still struggle with real data projects because of limited context windows, implicit schema relationships, ambiguous domain terminology, and the lack of reusable analysis experience. To address these problems, this thesis designs and implements Pontis, a graph-native agent workspace for project data. Instead of training a specialized Text-to-SQL model, Pontis introduces a queryable, updatable, and reusable knowledge graph layer between raw data and LLM agents, enabling agents to understand a data project before answering natural language queries.

This thesis first presents a workspace architecture centered on Neo4j and Cypher. Files, databases, tables, columns, foreign keys, statistical summaries, semantic descriptions, disambiguation hints, and cross-project experience are modeled as graph entities and relationships. Source modules publish source-derived facts such as file-system and SQLite schema information before graph queries. Extractor modules generate derived knowledge including column statistics, value distributions, column overlaps, foreign-key validation, semantic summaries, and vector indexes. All upper-level tools and agents access the graph through a unified \texttt{Workspace.cypher} interface. The thesis then presents a tool-calling data agent equipped with tools such as \texttt{find}, \texttt{meta}, \texttt{query}, \texttt{grep}, \texttt{read}, and \texttt{cypher}. Guardrails including round limitation, exploration checking, SQL entity checking, bridge-table checking, and SQL disambiguation checking are introduced to reduce hallucination, incorrect joins, and unnecessary exploration.

Preliminary experiments on the BIRD Text-to-SQL development set show that Pontis achieves 64.45\% execution accuracy on 1,533 logged questions. Error analysis indicates that the major failure sources are benchmark-specific SQL style preferences, local schema misunderstanding, value grounding errors, and join-path mistakes. Since real Text-to-SQL tasks often involve ambiguous natural language, underspecified user intent, and noisy reference SQL annotations, this thesis evaluates the system with both execution accuracy and system-level cost metrics. The planned full evaluation compares Pontis with Bash Agent, CHESS, Alpha-SQL, and DeepEye-SQL under a shared model configuration and private data copies. The comparison focuses on execution accuracy, cacheable pre-input tokens, runtime input tokens, runtime output tokens, LLM rounds, and total token consumption. The value of Pontis is not to replace specialized SQL models, but to provide a transparent, cacheable, and continuously evolving data-understanding workspace for general-purpose data agents.

\end{abstracten}

\keywordsen{large language model; agent; knowledge graph; Text-to-SQL; data workspace}

\tableofcontents

% ==================== 正文 ====================
\mainmatter

\section{绪论}

\subsection{研究背景与意义}

数据正在成为组织运行、科研发现和软件系统决策的核心基础设施。与早期主要围绕单一关系数据库或固定报表展开的数据访问不同，现代数据工作流通常横跨数据仓库、湖仓、业务语义层、BI 仪表盘、数据管道、实验记录、代码仓库和团队知识库。数据消费者也不再只满足于执行预定义 SQL，而是希望用自然语言提出探索性问题、追问分析过程、验证指标口径，并把得到的发现转化为后续决策或自动化动作。因此，数据系统面临的关键问题正在从“如何高效存储和查询数据”扩展为“如何让智能体理解、检索、验证并持续维护一个组织的数据知识”。

在这一背景下，数据智能体（data agent）成为学术界和工业界共同关注的方向。相比单次输入输出 Text-to-SQL 模型，数据智能体是能够围绕数据任务进行规划、调用工具、读取元数据、执行查询、验证结果并生成分析结论的系统。学术界的演进也体现了这一变化：Spider\upcite{yu2018spider} 和 BIRD\upcite{li2023bird} 推动了 Text-to-SQL 从跨库泛化走向真实数据库和执行约束；Spider 2.0\upcite{lei2025spider2} 进一步把评测对象扩展到真实数据应用和企业级工作流；InsightBench\upcite{sahu2025insightbench} 与 Data Agent Benchmark\upcite{dab2026dataAgentBench} 则将问题推进到多步洞察生成、异构数据源和企业数据问答。由此可见，自然语言数据分析正在从“单条 SQL 是否生成正确”转向“智能体能否在复杂数据环境中稳定完成分析任务”。

工业界也在快速把数据智能体嵌入主流数据平台。Microsoft Fabric data agent 被定位为 Fabric 中的对话式分析组件，能够连接 OneLake 中的 lakehouse、warehouse、semantic model 和 KQL 数据库，并在权限与治理约束下选择合适的数据源和查询工具\upcite{microsoft2026fabricDataAgent}。Snowflake Cortex Agents 通过 Cortex Analyst 和 Cortex Search 在结构化与非结构化数据之间进行工具路由，使应用可以在企业数据治理范围内完成 SQL 分析和文档检索\upcite{snowflake2026cortexAgents}。Databricks Genie Space 则要求领域专家用数据集、示例查询和文本指南配置面向业务团队的自然语言分析空间，并允许随着数据和用户问题变化持续更新语义知识\upcite{databricks2026genieSpace}。Tableau Next 和 agentic analytics 的提出也反映了 BI 平台从可视化报表向“从数据到洞察再到行动”的智能体化分析体验转变\upcite{tableau2025agenticAnalytics}。

这些趋势表明，真实的数据智能体系统必须同时处理四类能力：第一，能够在复杂数据资产中定位相关实体和上下文；第二，能够理解组织内部的业务术语、指标口径和历史经验；第三，能够通过工具在真实数据上验证假设，而不是只依赖模型参数中的常识；第四，能够在安全、权限和治理边界内运行，并为后续问题积累可复用知识。换言之，数据智能体的核心难点不只是 SQL 语法生成，而是上下文组织、语义 grounding、工具编排、评测可靠性和长期记忆的系统工程问题。企业内部系统的实践研究也显示，有效的 Text-to-SQL 系统需要索引数据库元数据、历史查询、文档和代码等多源知识，而不能只依赖表名和列名等简单元数据\upcite{chen2025enterpriseTextToSQL,johnson2024querygpt}。

因此，直接将大语言模型用于真实数据项目仍面临以下困难。

第一，数据库模式和项目上下文过大。真实数据库可能包含大量表、列、视图和项目文档。若将完整 schema、样例值和说明文档全部纳入提示词，会造成上下文冗长、成本上升和注意力分散；若预先过滤 schema，又可能错误丢弃关键表列。尽管如今随着大模型上下文窗口不断增大，数据分析任务可在数据库规模较小时直接全量输入数据库schema；但当企业 schema 继续扩张时，系统仍需要可靠的上下文控制机制来降低模型认知复杂度\upcite{maamari2024deathSchemaLinking}。

第二，关系、语义和值口径常常是隐式的。许多数据库没有完整外键约束，表间连接依赖命名约定、值域重叠、历史查询或项目说明。列名相同并不意味着语义相同，列名不同也可能表达相同业务概念；即使表列选择正确，过滤值、枚举别名、日期边界和聚合口径也仍可能出错。相关研究表明，列统计、值分布、查询日志、业务规则和 LLM 生成的字段描述可以显著补足人工文档缺失的问题\upcite{shkapenyuk2025metadata}\upcite{wang2025schemaGraphQA}\upcite{rissaki2024schemaRefinement}。

第三，模型必须在真实数据上验证假设。自然语言问题中的过滤值、枚举项、日期边界、比例分母和聚合粒度往往无法仅靠 schema 推断。SQL-PaLM、CHESS、ReFoRCE 和 APEX-SQL 等系统分别通过执行反馈、候选选择、多智能体协作和列探索来提升可靠性\upcite{sun2024sqlpalm}\upcite{talaei2024chess}\upcite{deng2025reforce}\upcite{cao2026apexsql}。这些工作说明，Text-to-SQL 正在从一次性生成转向带工具调用、检索、验证和修正的 test-time framework。

第四，评测结果本身需要谨慎解释。BIRD 等基准中存在自然语言歧义、gold SQL 口径差异和标注错误；已有研究发现，噪声和标注问题会影响模型排名和方法判断\upcite{wretblad2024birdNoise,jin2026annotationErrors}。这意味着系统评估不能只报告最终执行准确率，还应分析错误来源、可缓存上下文、运行时 token 成本和 LLM 调用轮次。

本文认为，真实数据项目中的自然语言分析不应只被视为单次 SQL 生成问题，而应被视为“智能体如何持续理解并维护一个数据项目”的问题。Pontis 因此在原始数据和 LLM 智能体之间引入图谱化工作空间，将 schema、统计、关系、摘要、消歧和经验统一表示为可查询知识，使智能体能够通过工具逐步获取精确上下文，而不是每道题从零扫描整个项目。

\subsection{主要贡献}

围绕上述问题，本文设计并实现了 Pontis 图谱化数据工作空间。本文的主要贡献如下：

（1）提出一种图谱化数据工作空间架构。该架构以 Neo4j 为持久化图存储，以 Cypher 为统一查询接口，以自动化脚本和智能体探索共同维护项目知识，避免在每次问答中直接拼接全部数据上下文。

（2）实现一套面向数据库项目的知识提取与组织机制。系统能够自动暴露文件系统、SQLite 模式、表列结构和外键关系，并进一步生成列级统计、值样例、列重叠、AI 摘要和语义向量等派生知识。

（3）设计了结合工具调用和护栏的 LLM 数据智能体。系统将图谱检索、元数据查看、SQL 执行、文件搜索和图写入操作封装为工具，并通过 SQL 实体检查、桥接表检查和消歧检查减少常见数据库问答错误。

（4）对当前方法与现有方法进行精度、效率等指标上的比较，并进行错因分析。本文将 Pontis 与 Bash Agent、CHESS、Alpha-SQL、DeepEye-SQL 等方法放在统一评测协议下比较，除准确率外统计 token 消耗、LLM 调用轮次等效率指标，使系统型贡献能够被量化分析。

\subsection{论文结构}

本文结构如下：第二章介绍相关工作；第三章给出 Pontis 的系统架构；第四章阐述自动化图谱构建与智能体知识探索机制；第五章介绍智能体工具系统与护栏机制；第六章给出实验设计、阶段性结果和指标统计；第七章总结全文并展望未来工作。

\section{相关工作}

\subsection{Text-to-SQL 基准与真实场景演进}

Text-to-SQL 是自然语言处理与数据库系统交叉领域的重要任务，其目标是将自然语言问题转换为可执行 SQL。早期方法多依赖规则和语义解析，随后 Seq2SQL\upcite{zhong2017seq2sql} 等方法将神经网络和强化学习引入 SQL 生成。Spider\upcite{yu2018spider} 的提出推动了跨数据库泛化研究，RAT-SQL\upcite{wang2020rat} 通过关系感知编码器建模问题与数据库 schema 之间的关系，代表了在深度模型时代显式建模 schema linking 的典型路线。

大语言模型出现后，Text-to-SQL 的研究重心从模型结构设计逐渐转向提示工程、检索增强、任务分解和执行反馈。C3\upcite{gao2024c3} 和 DAIL-SQL\upcite{gao2023dailsql} 系统比较了提示表示、示例选择和自修正等因素，DIN-SQL\upcite{pourreza2024dinsql} 则将复杂 SQL 生成拆分为 schema linking、查询分类和 SQL 生成等子任务。这些工作表明，强 LLM 不一定需要任务级微调才能完成复杂 SQL，但其效果高度依赖上下文组织和推理流程设计。BIRD\upcite{li2023bird} 进一步引入更大的数据库、外部 evidence、真实值理解和 SQL 效率要求，使 Text-to-SQL 从语法映射转向数据库内容理解。

近年的基准进一步强调企业级复杂性和端到端数据分析能力：一类工作将任务扩展到真实数据应用、企业级 schema、查询日志和 SQL workflow，另一类工作则把评测对象从单条 SQL 扩展到多步业务洞察生成\upcite{lei2025spider2}\upcite{chen2025beaver}\upcite{sahu2025insightbench}。这些工作共同说明，真实 Text-to-SQL 系统需要处理远大于教学数据集的 schema 空间、业务术语、私有规则、SQL 方言和多轮分析流程。本文关注的 Pontis 沿着这一趋势，希望把 Text-to-SQL 问题从简单生成任务拓展到“项目级数据理解”。

\subsection{Schema Linking 与元数据增强}

Schema linking 负责将自然语言问题中的实体、属性和约束映射到数据库表列，是 Text-to-SQL 系统的核心环节。传统方法通常将 schema linking 作为 SQL 生成前的过滤步骤，但近期研究对这一环节提出了更细的分析。Maamari 等人的研究指出，当完整 schema 能放入上下文且模型推理能力足够强时，过早过滤 schema 可能因为召回错误而伤害生成结果\upcite{maamari2024deathSchemaLinking}。这一结论说明，schema linking 不能被简单理解为“越早过滤越好”，其价值取决于 schema 规模、上下文窗口、模型能力和召回风险。对 Pontis 而言，关键不只是裁剪 schema，而是把可复用的上下文组织为可查询的项目知识，使智能体在需要时逐步获取证据。

单纯依赖表名和列名也不足以解决真实场景中的语义问题。自动元数据抽取研究从数据库 profiling、列级统计、值样例、历史查询日志和 LLM 生成字段摘要中提取补充信息\upcite{shkapenyuk2025metadata}；schema refinement 工作则把原始数据库整理为更易解释的语义层，降低用户问题与底层 schema 之间的语义距离\upcite{rissaki2024schemaRefinement}。面向多表问答的 schema graph 方法进一步说明，表列之间的连接关系、候选路径和人工确认信息适合用图结构表达，而不应完全依赖模型在生成时临场推断\upcite{wang2025schemaGraphQA}。

除此之外，schema linking 只解决“选择哪些表列”的问题，并不能自动保证过滤值、枚举映射、日期边界和聚合口径正确。BIRD 的设计已经强调数据库值和外部 evidence 的作用\upcite{li2023bird}，而企业系统中常见的业务简称、历史口径和隐式指标定义则使 value grounding 更加困难。RubikSQL 将 NL2SQL 视为长期学习任务，持续维护业务术语、规则和失败经验\upcite{chen2025rubiksql}；LinkedIn 的企业 Text-to-SQL 系统也把元数据、查询日志、文档和代码共同纳入知识图谱\upcite{chen2025enterpriseTextToSQL}。这些工作表明，真实数据项目中的上下文不只是 schema 文本，而是 schema、值、关系和经验的组合。

Pontis 与这些工作的共同点在于重视元数据、值分布和关系知识；不同点在于，Pontis 将这些信息统一落到项目级图谱工作空间中，并通过工具接口让智能体按需读取，而不是一次性把全部增强 schema 拼入 prompt。

\subsection{智能体式 Text-to-SQL 与执行验证}

工具调用是增强 LLM 能力的重要途径。在 Text-to-SQL 中，工具调用可以弥补模型静态知识不足的问题：模型可以执行 SQL 验证假设，通过搜索定位字段，通过样例值确认枚举，通过错误反馈修正语法和逻辑。

围绕 BIRD 和 Spider 2.0，已有方法逐渐形成多阶段、可验证的 test-time framework。SQL-PaLM\upcite{sun2024sqlpalm} 研究了 few-shot prompting、执行过滤、consistency decoding、指令微调和数据库内容使用对性能的影响；CHESS\upcite{talaei2024chess} 使用信息检索、schema 选择、候选生成和单元测试等多个 agent 分工协作；ReFoRCE\upcite{deng2025reforce} 引入表压缩、格式约束、自修正、共识选择和列探索；APEX-SQL\upcite{cao2026apexsql} 将生成过程组织为假设验证循环，通过真实数据探索落地语义。上述工作说明，复杂 Text-to-SQL 的关键并不只是更强的单次生成，而是如何组织检索、探索、验证和选择。

除显式工具调用外，候选集选择和 test-time scaling 也是提升稳定性的常见路径。self-consistency\upcite{wang2023selfConsistency} 说明多条推理路径的一致性可以提高复杂推理可靠性，Text-to-SQL 系统也常借助多候选生成、执行筛选和一致性选择来降低偶然错误。这类方法提升了复杂任务上的准确率，但通常也会增加 token 消耗和 LLM 调用轮次，因此需要和统一实验指标一起分析。

Pontis 也采用工具智能体范式，但其重点不是构造更多候选 SQL，而是提供一个可持久化的数据理解层。智能体的每次检索、元数据查看和 SQL 验证都围绕图实体展开，后续问题可以复用已有的统计、摘要、消歧和经验知识。

\subsection{知识图谱与数据项目理解}

知识图谱将实体及其关系表示为图结构，适合表达数据库 schema、文件组织、表列归属、外键关系和业务概念之间的联系。传统知识图谱构建强调从文本中抽取实体和关系，而数据项目场景中的图谱构建还需要处理结构化源事实和统计派生知识。企业 Text-to-SQL 实践已经显示，知识图谱可以把数据库元数据、历史查询、文档和代码等多源信息组织为可检索上下文，为 SQL 生成和错误修复提供依据\upcite{chen2025enterpriseTextToSQL}。

在企业数据分析场景中，知识图谱的价值不仅在于存储概念关系，还在于为智能体提供可操作的项目地图。Schema graph 相关工作使用图结构减少模型对隐式连接关系的猜测\upcite{wang2025schemaGraphQA}；RubikSQL 将长期运行中沉淀的业务知识、错误修正和查询经验组织为可复用知识库\upcite{chen2025rubiksql}；LinkedIn 的企业系统使用知识图谱统一数据资产、元数据和查询经验\upcite{chen2025enterpriseTextToSQL}；Uber QueryGPT 则结合向量数据库和相似查询检索，将自然语言问题和内部数据模型知识连接起来\upcite{johnson2024querygpt}。这些实践说明，Text-to-SQL 系统需要一种位于数据库和模型之间的中间表示，用来承接不断变化的项目知识。

Pontis 的图谱并不是静态业务知识库，而是围绕项目数据持续生成和更新的工作空间。文件系统和数据库 schema 可以作为源事实按需刷新，AI 摘要、列统计、列重叠和人工经验可以作为派生事实持久化。通过这种方式，Pontis 将“当前项目是什么”和“过去分析学到了什么”统一到一个可查询图中。与只服务于单次 SQL 生成的增强 prompt 相比，图谱工作空间更适合跨问题复用、错误归因和后续知识更新。

\subsection{评测噪声与系统指标}

Text-to-SQL 评测通常依赖执行准确率，即比较预测 SQL 与参考 SQL 的执行结果是否一致。然而，执行结果并不总能完整反映语义正确性：不同 SQL 可能因输出列、排序、聚合口径或边界条件差异而产生不同结果；模型可能为了可读性等因素生成了优于参考sql的结果，但却因为格式问题被判错；参考 SQL 本身也可能存在标注错误或问题歧义。针对 BIRD 的噪声分析指出，问题和 gold SQL 中的错误会影响模型表现和结论可靠性\upcite{wretblad2024birdNoise}；进一步的 benchmark 审计研究也表明，标注错误可能显著改变模型排名\upcite{jin2026annotationErrors}。

随着基准复杂度提高，评测对象也从“单条 SQL 是否匹配”扩展为“系统是否能稳定完成分析任务”。因此，系统比较需要同时报告错误类型和资源消耗，尤其要区分 schema linking 失败、value grounding 失败、SQL 语义错误、输出口径差异和评测噪声。

基于这一考虑，本文在评估 Pontis 时不仅关注执行准确率，还关注预输入 token、runtime 输入 token、runtime 输出 token、LLM 调用轮次和总 token 消耗。这些指标更适合衡量图谱化工作空间的价值：即它是否能把可缓存的稳定上下文和每题运行时上下文区分开，减少重复输入，并提高 schema/value grounding 的可靠性。

\textcolor{red}{此外，真实数据分析系统中的查询通常不是完全独立同分布的样本。用户会围绕同一数据库持续提出相关问题，历史正确 SQL、失败案例、字段消歧和业务口径都可能影响后续回答。若评测强制把同一数据库下的每个问题视为相互独立的单次生成任务，就会低估经验复用和增量知识维护的价值。因此，本文在保留传统逐题执行准确率的同时，也将 Pontis视为能够随项目查询历史持续更新的工作空间，后续实验可进一步比较独立评测与 query-log 式连续评测之间的差异。}
\textcolor{red}{【有待商榷，尽管企业确实存在通过上一条query的golden sql（或者通过企业已有的sql log历史）来改进下一条query,进行经验沉淀的需求，但目前没有text2sql方法采用这种非独立同分布的实验逻辑】}


\section{Pontis 系统架构}

\subsection{总体设计}

Pontis 的核心思想是在原始数据源和智能体之间建立图谱化中间层。系统整体由四个部分组成：原始数据源、图谱工作空间、知识提取器和工具智能体。

Pontis 首先通过确定性的自动化脚本将这些数据源中的基础事实发布为图节点和边，例如目录、文件、数据库、表、列和外键；随后继续在图谱之上生成统计、摘要、关系和向量等增强信息。Neo4j 图谱工作空间负责持久化这些事实和派生知识，并提供 Cypher 查询接口。工具智能体通过工具访问图谱和原始数据，逐步完成自然语言分析任务。

系统的核心接口边界是 \texttt{Workspace.cypher(...)}。上层工具和知识提取流程均通过该接口访问图谱，避免直接依赖底层图存储的具体实现。该设计将数据源刷新、图查询和结果指针解析逻辑集中在工作空间层，从而降低工具层与存储层之间的耦合。

\subsection{图谱数据模型}

Pontis 将项目数据表示为属性图。在典型的数据库Text-to-SQL场景中，主要涉及的实体类型包括：

（1）文件实体。表示项目目录中的文件，包含路径、文件类型和源访问指针等属性。

（2）数据库实体。表示 SQLite 等数据库文件，作为表、视图和列的上级节点。

（3）表与视图实体。表示数据库表、视图或表格文件中的逻辑表，记录名称、来源和结构信息。

（4）列实体。表示表中的字段，记录列名、类型、样例值、统计信息和语义描述。

（5）外键和关系实体。表示显式外键、推断关系、列重叠和消歧知识。

（6）知识实体。表示人工或自动沉淀的规则、经验、错误分析结果和 benchmark 案例。


\subsection{Neo4j 与 Cypher 接口}

Pontis 使用 Neo4j 作为图存储后端。工作空间层对外暴露 Cypher 接口，基础事实发布脚本和工具通过 Cypher 语句完成图实体写入、匹配和查询。选择 Neo4j 的主要原因包括：

（1）Cypher 能直接表达图模式匹配、邻居遍历和多跳关系查询，适合表列关系、外键关系和知识实体检索。

（2）图存储与工具层解耦。工具只需要关心实体 ref、标签和查询结果，不需要自己维护图索引。

（3）跨项目路由更清晰。Pontis 可以为不同项目配置不同 Neo4j 端口或数据库，使同一套工具能够访问多个项目图。

（4）后续可利用 Neo4j 的索引、向量检索和查询优化能力，支撑更大规模数据项目。

\section{图谱构建与知识提取机制}

第四章介绍 Pontis 的图谱构建与知识提取机制。面向真实数据项目的自然语言分析任务，系统需要处理的上下文并不限于数据库 schema。除表名、列名和外键等结构化事实外，列值分布、字段语义、隐式连接关系、业务术语、历史经验和非结构化说明文档都会影响最终 SQL 的正确性。传统 Text-to-SQL 方法通常在推理时临时拼接 schema 文本，而企业数据和真实 benchmark 中常见的弱文档、弱约束和多源数据使这种方式难以稳定覆盖必要证据。因此，Pontis 将项目知识组织为可增量更新的图谱工作空间，使工具智能体能够以实体和关系为单位读取上下文，而不必在每次问答时重新扫描完整项目。

现有研究将空值率、不同值数量、数据类型、值模式、函数依赖和包含依赖等信息视为理解数据集的重要元数据\upcite{abedjan2015profiling}；Aurum 等数据发现系统通过图结构组织企业数据资产及其关系，以支持跨数据源发现和查询\upcite{fernandez2018aurum}；自动元数据抽取和 schema graph 相关工作则进一步说明，列统计、字段描述、连接路径和人工确认知识能够补足原始 schema 的语义缺口\upcite{shkapenyuk2025metadata}\upcite{wang2025schemaGraphQA}。Pontis 的目标不是单独复现其中某一种技术，而是把这些知识类型纳入同一运行时工作空间，为后续的检索、探索、验证和经验复用提供基础。

Pontis 的图谱构建采用分阶段、可增量的组织方式。首先，自动化脚本把文件系统、文本文件、表格文件和数据库模式转换为图节点与边，形成稳定的源事实视图；其次，知识提取流程在源事实之上补充列统计、值分布、结构模式、关系候选、语义摘要和向量索引；最后，智能体探索流程针对难以完全规则化的语义关系和消歧问题生成经验性知识。三类结果共享同一图谱存储，但在来源、生成方式和可信度上保持区分。确定性事实和统计画像主要来自脚本，语义判断和连接经验主要来自智能体探索，从而降低错误推断对基础事实的污染风险。

\subsection{节点身份、边与访问指针}

Pontis 在图谱中同时维护“实体身份”和“原始数据访问能力”。实体身份由路径、引用和内部 id 共同保证：文件节点通常以路径合并，表和列节点使用数据库路径、表名和列名组成稳定引用，外键和关系实体则通过两端表列确定。所有节点都会获得存储层生成的内部 id，用于稳定边连接和工具显示。稳定身份是图谱复用的前提：如果同一张表、同一列或同一段文本在不同运行中无法被识别为同一实体，统计结果、摘要、向量和智能体经验就难以增量更新，也无法可靠地被后续问题复用。

图谱建模遵循一个原则：具有语义的内容均建模为节点，例如数据库中无论是表、列、外键还是业务术语、历史经验、错误分析结果，都以节点形式存在；而边不附加类型和属性，仅用于表征图谱间实体的关系性，便于agent直接探索和查询。

从数据发现角度看，图谱中的边主要作为被智能体利用的导航路径。文件到文本 chunk、数据库到表、表到列、列到统计、列到候选关系、实体到摘要和向量索引之间的连接，使系统能够从一个自然语言查询逐步缩小到相关数据源、相关表列和可能的连接路径。Schema graph 工作表明，显式维护表列之间的图结构能够减少模型在多表问答中对隐式关系的临场猜测\upcite{wang2025schemaGraphQA}。Pontis 将这一思想扩展到项目级数据资产，使关系型表、半结构化文件和智能体生成的经验知识可以处在同一个查询空间中。


\subsection{自动化脚本式知识提取}

自动化脚本式知识提取负责生成可以由数据源稳定推导或批量计算的图谱知识。每个提取单元接收一个 \texttt{Workspace} 对象，通过 Cypher 和工具函数读取源事实，并将结果写回图谱。其任务既包括提取文件、文本、表格、数据库、表、列和外键等基础事实，也包括在基础事实之上生成统计、摘要、关系候选和向量索引。因此，这类脚本可以离线批量运行，也可以按实验需要选择性执行。

该部分的核心价值在于把大量低层、重复、可计算的观察过程从在线智能体循环中前移。对于一个新项目，智能体如果在回答每个问题时都临时查看目录、浏览 schema、执行 \texttt{SELECT DISTINCT}、采样字段值并判断候选关系，会造成较高的 LLM 轮次和 token 消耗。自动化提取将这些步骤转化为可缓存的图谱事实，使在线阶段更多地围绕已提取证据进行推理。数据画像和自动元数据抽取研究均表明，这类统计性和描述性元数据能够显著改善数据理解过程\upcite{abedjan2015profiling}\upcite{shkapenyuk2025metadata}。

系统中的主要提取单元如表\ref{tab:extractors} 所示。

\begin{table}[htbp]\small
  \centering
  \topcaption{Pontis 主要知识提取单元}
  \label{tab:extractors}
  \begin{tabularx}{\textwidth}{lX}
    \toprule[1.5pt]
    \textbf{提取单元} & \textbf{功能} \\
    \midrule[1pt]
    文件系统提取单元 & 提取目录、文件实体和文本文件元信息 \\
    数据库结构提取单元 & 提取 SQLite 数据库、表、列和外键 \\
    数据库列统计模块 & 提取 SQLite 列统计、样例值和高频值 \\
    外键验证模块 & 验证数据库声明外键的有效性 \\
    列值重叠发现模块 & 计算列值重叠并生成潜在连接关系候选 \\
    列语义摘要模块 & 使用 LLM 生成数据库列的语义摘要 \\
    语义向量模块 & 为实体的 detail 文本生成语义向量 \\
    局部消歧模块 & 生成表列和业务术语的局部消歧知识 \\
    智能体连接发现模块 & 通过智能体探索潜在连接关系和 join 经验 \\
    \bottomrule[1.5pt]
  \end{tabularx}
\end{table}

自动化提取框架具有三个设计特征。第一，提取单元粒度较细，可以根据实验需要选择只运行源事实、统计、摘要或语义向量模块。第二，执行顺序可配置，既可以先运行基础结构提取和轻量统计，再运行关系发现与 LLM 摘要，也可以在 benchmark 中跳过高成本模块。第三，提取结果回写到同一个图谱工作空间中，后续工具和智能体无需感知提取过程的内部实现，只需查询图实体。

在典型数据库项目中，自动化流水线可以分为四个阶段。第一阶段读取表、列和外键等基础节点，建立基础图；第二阶段扫描真实数据，生成列统计、样例值和高频值；第三阶段根据外键、值域重叠和统计特征生成关系候选；第四阶段使用 LLM 或 embedding 服务生成语义摘要和向量索引。各阶段的输出都以节点属性或派生实体的形式写回图谱，因此后续阶段可以复用前序阶段的结果。

\subsection{列统计与值分布}

列统计是 Text-to-SQL 中的重要上下文。仅靠列名往往无法判断一个字段的真实含义，例如 \texttt{type}、\texttt{id}、\texttt{status} 等列名在不同表中可能含义不同；相反，字段中的真实取值、空值比例、基数和分布形态往往能揭示其业务角色。数据画像研究将这类信息归入单列元数据，并进一步讨论唯一列组合、函数依赖和包含依赖等跨列特征\upcite{abedjan2015profiling}。Pontis 在系统实现上先聚焦对在线问答最直接有用的轻量统计，对数据库和表格文件中的列提取以下信息：

（1）基本类型信息，包括 SQLite 声明类型和归一化类型。

（2）空值、不同值数量、样例值和文本长度等轻量统计。

（3）数值列的范围和分布摘要。

（4）文本列的高频值，用于帮助模型理解枚举字段和业务取值。

这些统计信息会写入列实体的元数据，并在 \texttt{meta} 工具中呈现。与直接把整列数据交给 LLM 不同，统计摘要将大量原始值压缩成稳定、低成本的上下文。对于大列，系统不要求保存完整 distinct 集合，而是使用近似基数估计和流式高频项统计来降低预处理成本；对于文本列，系统保留少量去重样例和 top-k 高频值，使智能体能够判断枚举值、编码格式和自然语言问题中的值映射。

列统计同时服务于性能和准确率。性能上，它避免智能体为了理解一个枚举列而反复执行 \texttt{SELECT DISTINCT}；准确率上，它帮助模型判断过滤条件是否需要大小写转换、是否存在缩写、是否应该使用编码列还是描述列。许多 value grounding 错误并非源于表列召回失败，而是源于模型缺乏列中真实取值的证据。APEX-SQL 也强调通过真实数据探索验证列角色和查询假设\upcite{cao2026apexsql}。Pontis 将这类证据预先压缩到图谱中，使在线智能体可以用一次元数据读取获得候选值、分布和异常线索，从而降低 SQL 探索轮次。

列统计还为后续语义摘要和关系发现提供输入。列摘要模块需要依赖样例值、高频值和基数判断字段是否为枚举、标识符、金额、日期或自然语言描述；列值重叠模块则需要利用 distinct 值和覆盖率判断两列是否可能构成连接路径。因此，统计画像并不是孤立的辅助信息，而是 Pontis 图谱中从原始数据到高层语义知识的中间层。

\subsection{智能体知识探索}

自动化脚本适合生成稳定事实和统计画像，但真实数据项目中仍有一类知识难以完全规则化：字段语义需要结合上下文解释，隐式连接关系需要在全局 schema 中验证，同名列和近义字段需要根据样例值区分，长文档和说明材料中的业务口径也需要被压缩成可执行规则。Pontis 将承担这类任务的预处理智能体称为 explorer。与最终回答用户问题的 query agent 不同，explorer 的目标不是直接生成 SQL，而是在离线或半离线阶段把可复用的数据理解写回图谱。

Explorer 的设计并不是把一个通用 agent 放到项目中自由探索，而是采用专门化任务、受限工具和保守写入的组合。Toolformer 和 ReAct 说明了语言模型可以通过工具调用在外部环境中获取信息并逐步推理\upcite{schick2023toolformer}\upcite{yao2023react}；CHESS、ReFoRCE 和 APEX-SQL 进一步把 Text-to-SQL 拆解为信息检索、schema 选择、列探索、候选生成和验证等子过程\upcite{talaei2024chess}\upcite{deng2025reforce}\upcite{cao2026apexsql}。Pontis 的 explorer 采用类似的任务分解思想，但输出对象不是一次性 SQL 候选，而是图谱中的长期知识。系统为 explorer 显式指定 prompt、工具集合和写入范围，通常只启用与任务相关的 \texttt{find}、\texttt{meta}、\texttt{query}、\texttt{update\_meta}、\texttt{create\_entity} 和 \texttt{add\_edge} 等工具，并通过轮次护栏限制无界探索。

系统中的主要 explorer 如表\ref{tab:explorers} 所示。它们按照知识类型分工，而不是按照代码目录简单分类：摘要类 explorer 负责把表、列、文件和结构模式转化为自然语言语义；关系类 explorer 负责把候选连接线索提升为高置信度关系；消歧类 explorer 负责记录容易误用的字段差异。

\begin{table}[htbp]\small
  \centering
  \topcaption{Pontis 主要 explorer 的设计职责}
  \label{tab:explorers}
  \begin{tabularx}{\textwidth}{lXX}
    \toprule[1.5pt]
    \textbf{Explorer} & \textbf{设计思路} & \textbf{写入的知识} \\
    \midrule[1pt]
    \texttt{analyze} & 依赖感知的全库理解。先读 README 和数据库入口，再按核心表、依赖表、关系实体逐步展开；复杂表可交给子智能体，避免单轮上下文过大。 & 表、列、数据库、关系实体的 \texttt{brief/detail}，特别是字段业务含义、值特征和易混淆说明。 \\
    \texttt{join\_detect} & 高置信度关系发现。以外键和 \texttt{overlap} 为候选线索，结合列统计、样例值和全局一致性检查，只在证据充分时创建关系。 & \texttt{rel} 实体及其证据说明、使用注意和不确定性声明。 \\
    \texttt{disambiguate} & 语义消歧。扫描同名、近名或同义不同名的表列，要求基于真实数据和现有关系判断差异，而不是只看字段名。 & \texttt{disambig} 实体，以及相关列的补充 \texttt{detail}。 \\
    \texttt{readme} & 项目级知识压缩。复用已生成的表列摘要、关系和消歧结果，形成面向后续 agent 的库级导航文档。 & README 节点的项目概览、主要对象、关系结构、数据质量风险和建议探索路径。 \\
    \texttt{text\_chunk} & 长文本结构化。先理解文档主题，再按语义片段创建 chunk；长文件由子智能体分段阅读，主智能体统一写图。 & 文本文件摘要、chunk 实体及其行号范围和主题说明。 \\
    \bottomrule[1.5pt]
  \end{tabularx}
\end{table}

其中，\texttt{join\_detect} 体现了 Pontis 对关系知识的保守态度。真实数据库中并非所有连接关系都有显式外键，尤其在 benchmark 数据库、历史业务库和数据湖表格中，连接规则常常只存在于命名约定、值域重叠或业务经验中。数据发现系统通常需要在大量数据资产之间发现相关表、可连接列和可联合数据集\upcite{fernandez2018aurum}\upcite{fan2023starmie}；schema graph 方法也说明，多表问答需要显式维护可解释的连接路径，而不能完全依赖 LLM 在生成时猜测\upcite{wang2025schemaGraphQA}。Pontis 因此同时利用三类信息发现关系：数据库声明外键、列值重叠候选和智能体验证后的 \texttt{rel} 实体。列值重叠模块批量生成候选证据，\texttt{join\_detect} 则读取外键、overlap、列统计和样例值，检查是否已有确定关系覆盖、是否存在传递冗余、是否与全局结构冲突，只有在高置信度下才写入 \texttt{rel}。这种“候选生成 + 保守确认”的设计接近 APEX-SQL 中用真实数据探索验证假设的思路\upcite{cao2026apexsql}，但 Pontis 将验证结果固化为可复用图谱知识。

\texttt{disambiguate} 处理的是另一类高频错误：字段名相同不代表语义相同，字段名不同也可能指向同一业务概念。例如多个表中都存在 \texttt{id}、\texttt{name} 或 \texttt{type} 时，最终 SQL 生成前需要确认字段来源；而 \texttt{Free Meal Count} 与 \texttt{FRPM Count} 这类近义字段也可能存在包含关系或统计口径差异。该 explorer 会按列名和近名模式收集候选实体，再通过列统计、样例值、表职责和已有关系判断是否真的存在歧义。确认后，系统创建 \texttt{disambig} 实体并补充相关列的 \texttt{detail}，后续 SQL 消歧护栏会在高风险输出时提醒 query agent 读取这些信息。Agentic schema refinement 强调通过语义层整理降低用户问题与底层 schema 的距离\upcite{rissaki2024schemaRefinement}；Pontis 的消歧 explorer 可以看作面向局部歧义的轻量 schema refinement。

\texttt{analyze} 和 \texttt{text\_chunk} 则负责把数据资产转化为可检索、可复用的语义描述。\texttt{analyze} 不是简单遍历所有实体，而是按 README、数据库、核心表、依赖表和关系实体的顺序建立理解，并优先读取精确属性，避免一次性展开全项目导致上下文膨胀；文本 explorer 则为长文档创建语义 chunk。这些设计对应于数据画像、语义类型检测和数据湖发现研究中的共同观点：数据理解需要结合结构、统计、值模式和上下文表示，而不是只依赖文件名或字段名\upcite{abedjan2015profiling}\upcite{hulsebos2019sherlock}\upcite{fan2023starmie}。

\texttt{readme} 体现了 Pontis 对长期知识的利用。其把已有摘要、关系、消歧和数据质量线索压缩成库级 README，使后续 agent 第一次进入项目时不必从零开始。RubikSQL 将工业 NL2SQL 视为长期学习任务，强调通过数据库画像、规则挖掘和知识库维护持续积累领域知识\upcite{chen2025rubiksql}；LinkedIn 和 Uber 的企业实践也显示，文档、业务规则和知识图谱能够改善自然语言到 SQL 的稳定性\upcite{chen2025enterpriseTextToSQL}\upcite{johnson2024querygpt}。

\subsection{语义摘要与向量索引}

Pontis 中的 \texttt{brief} 和 \texttt{detail} 是面向智能体的核心文本属性。列统计提供字段取值层面的证据，AI 摘要进一步补充字段或表的业务语义。系统会对表、列和文本文件生成摘要，并通过语义向量模块为 detail 文本生成向量。语义类型检测和表理解研究表明，字段含义通常需要结合列名、值分布、字符模式和表上下文判断，而不能只依赖表头字符串\upcite{hulsebos2019sherlock}。因此，Pontis 的摘要生成也不是简单改写列名，而是把统计画像、样例值、上下游关系和文件上下文压缩成适合智能体读取的自然语言说明。

列摘要模块在生成列说明时输入列名、类型、基数、空值比例、数值范围、样例值和高频值，并要求输出稳定的业务语义，避免生成易过时的具体统计数字。对于表摘要，系统会聚合同表列信息和外键关系，描述表的职责、粒度和常见关联路径。对于文本文件，系统会提取文本分片摘要，使非结构化说明材料也能进入同一个检索空间。

语义向量模块只对存在非空 \texttt{detail} 的节点生成 embedding，并记录模型名、向量维度和 detail 文本 hash。当 detail 未变化且 embedding 配置一致时，模块会跳过该节点，避免重复请求 embedding 服务。向量写回 Neo4j 后，系统按节点标签创建向量索引，使 \texttt{find(ref, query)} 从名称匹配扩展为基于自然语言语义的实体检索。例如，当问题涉及“percentage”“ratio”“foreign player”或“school SAT score”时，系统可以召回对应表列、知识规则或历史案例。Starmie 等数据发现工作利用列表示和高效向量索引提升数据湖中的语义检索效率\upcite{fan2023starmie}；Pontis 在规模较小的项目知识图谱中采用相同方向的思想，使智能体更快聚焦到候选实体，减少全量浏览 schema 的成本。

\subsection{增量刷新与成本控制}

图谱构建的成本主要来自三部分：源扫描、数据统计和 LLM/embedding 调用。Pontis 对这三部分分别采用不同策略。源扫描通过数据源 fingerprint 和刷新缓存避免重复发布未变化的数据源；数据统计尽量使用单次扫描、近似基数和 top-k 结构压缩原始值；LLM 摘要和语义向量则通过缺失字段检查、文本 hash 和批量请求减少重复调用。

这种设计使 Pontis 能在预处理成本和在线查询成本之间进行明确取舍。预处理阶段预先投入计算成本生成统计、摘要和关系线索，可以减少智能体在线阶段的工具调用轮次、SQL 探索次数和输入 token。另一方面，并非所有项目都需要运行完整流水线：在轻量场景中可以只运行基础事实发布和列统计；在复杂 Text-to-SQL benchmark 中，则可以追加列重叠、AI 摘要、向量索引和智能体探索，以换取更强的 schema linking、value grounding 和 join-path 可靠性。

\section{智能体设计}

\subsection{智能体循环}

Pontis 智能体采用工具调用循环。每轮循环中，系统将用户问题、系统提示、工具定义和历史消息发送给 LLM；若 LLM 返回工具调用，Pontis 执行对应工具并把结果作为下一轮输入；若 LLM 返回最终文本，则暂时结束循环等待用户反馈。

该过程可以表示为：

\begin{equation}
A_t = \mathrm{LLM}(Q, H_{t-1}, T),
\end{equation}

其中 $Q$ 表示用户问题，$H_{t-1}$ 表示历史上下文和工具结果，$T$ 表示可用工具集合，$A_t$ 表示第 $t$ 轮模型动作。若 $A_t$ 是工具调用，则系统执行工具并更新 $H_t$；若 $A_t$ 是最终 SQL 或自然语言答案，则结束循环。

工具智能体的在线成本主要由预输入 token、runtime 输入 token、runtime 输出 token 和 LLM 调用轮次共同决定。预输入 token 对应可被 token 缓存命中的稳定前缀，runtime 输入 token 对应每题运行时新增上下文，runtime 输出 token 对应模型生成内容，LLM 调用轮次则反映串行模型交互次数。

\subsection{工具集合}

Pontis 将工具划分为只读工具、写入工具和子智能体工具。只读模式是 Text-to-SQL benchmark 的主要模式，包含实体查找、元数据读取、文件读取、内容搜索、SQL 查询和图查询等能力。写入模式增加知识实体创建、元数据更新、边添加和删除能力，用于经验沉淀和项目知识维护。

\begin{table}[htbp]\small
  \centering
  \topcaption{Pontis 主要工具}
  \label{tab:tools}
  \begin{tabular}{lll}
    \toprule[1.5pt]
    \textbf{工具} & \textbf{类别} & \textbf{作用} \\
    \midrule[1pt]
    \texttt{find} & 图检索 & 按 ref 模式和语义 query 查找实体 \\
    \texttt{meta} & 元数据 & 查看实体属性和邻居关系 \\
    \texttt{query} & 数据查询 & 在数据库或表格文件上执行只读 SQL \\
    \texttt{grep} & 文本搜索 & 在文件实体范围内搜索内容 \\
    \texttt{read} & 文件读取 & 读取文本文件指定行范围 \\
    \texttt{bash} & 命令执行 & 执行受限 shell 命令 \\
    \texttt{create\_entity} & 图写入 & 创建知识实体 \\
    \texttt{update\_meta} & 图写入 & 更新实体元数据 \\
    \texttt{add\_edge} & 图写入 & 添加实体关系边 \\
    \texttt{delete} & 图写入 & 删除实体 \\
    \texttt{agent} & 子任务 & 调用子智能体完成局部任务 \\
    \bottomrule[1.5pt]
  \end{tabular}
\end{table}

其中，\texttt{find}、\texttt{meta} 和 \texttt{query} 是 Text-to-SQL 场景中最核心的三个工具。\texttt{find} 用于缩小候选表列或知识范围，\texttt{meta} 用于精读实体语义和邻居关系，\texttt{query} 用于验证 SQL 假设和检查样例值。

\subsection{提示组装}

Pontis 的提示词并非单一固定 prompt，而是由稳定系统提示和任务实例提示共同构成。系统提示保留跨任务复用且对智能体理解数据项目必要的内容，包括基础身份、工具说明、图谱本体、SQL 规范、护栏规则、项目上下文和 README 摘要。

这种分层组装方式具有两个作用。第一，稳定系统提示可以在同一类任务中复用，减少重复 token 并有利于上下文缓存；第二，任务实例提示保留具体场景的约束，避免将评测输出格式、预处理写图策略或反思复盘规则混入所有智能体实例。工具相关提示词由创建智能体时显式传入的工具集合控制。

\subsection{护栏机制}

Pontis 的护栏层位于 LLM 输出和工具执行之间。每个护栏独立检查待执行工具调用或最终文本输出，并给出允许、警告或阻断裁决。多个护栏的裁决按阻断高于警告、警告高于允许的优先级聚合。

\begin{table}[htbp]\small
  \centering
  \topcaption{Pontis 护栏机制}
  \label{tab:guardrails}
  \begin{tabular}{lll}
    \toprule[1.5pt]
    \textbf{护栏} & \textbf{文件} & \textbf{作用} \\
    \midrule[1pt]
    RoundLimit & \texttt{round\_limit.py} & 达到轮次上限后阻止继续调用工具 \\
    ExplorationCheck & \texttt{exploration\_check.py} & 阻止过宽泛的全图探索 \\
    ToolAbuse & \texttt{tool\_abuse.py} & 限制指定工具的过量使用 \\
    SQLEntityCheck & \texttt{sql\_check.py} & 检查 SQL 引用表列是否已读且存在 \\
    BridgeTableCheck & \texttt{sql\_join\_check.py} & 检查多表连接是否遗漏桥接关系 \\
    SQLDisambigCheck & \texttt{sql\_disambig\_check.py} & 提醒同名列和局部术语消歧 \\
    \bottomrule[1.5pt]
  \end{tabular}
\end{table}

护栏的设计目标不是替代模型推理，而是在高风险节点提供结构化反馈。例如，模型若尝试输出最终 SQL，但其中引用的表列从未通过 \texttt{meta} 阅读，SQL 实体检查会阻断最终输出并要求模型先核实实体；若 join 中涉及可能需要桥接表的列对，桥接表检查会提示相关关系实体；若 SQL 中出现与消歧实体相关的术语，消歧检查会要求模型读取对应解释。

\section{实验设计与阶段性结果}

\subsection{研究问题}

本文实验围绕以下问题展开：

（1）Pontis 相比只给 shell/sqlite/file 工具的 Bash Agent，是否能通过图谱化数据工作空间提升 Text-to-SQL 准确率和稳定性？

（2）Pontis 在解题过程中产生的 LLM 调用轮次和 token 资源消耗具有什么特点？

（3）BIRD 等 Text-to-SQL 数据集中的错误样本主要来自哪些原因，系统错误与数据集口径或标注噪声之间应如何区分？

\subsection{数据集与评测指标}

本文采用 BIRD 开发集作为主要实验数据。BIRD 包含真实数据库、自然语言问题、evidence、标准 SQL 和难度标签，适合评估真实数据库问答能力。需要指出的是，BIRD 中部分题目具有较强 benchmark 风格，百分比、比例、输出列、排序和 \texttt{COUNT DISTINCT} 等口径可能存在多种业务上合理的解释。因此，本文把准确度作为必要但不充分的指标，并结合 token 资源消耗分析系统行为。

主准确度指标为执行准确率（Execution Accuracy, EX），即预测 SQL 与 gold SQL 在同一数据库上执行后的结果集合是否一致。除准确度外，本文还记录以下实验指标：

\begin{table}[htbp]\small
  \centering
  \topcaption{实验指标}
  \label{tab:metrics}
  \begin{tabularx}{\textwidth}{lX}
    \toprule[1.5pt]
    \textbf{指标} & \textbf{含义} \\
    \midrule[1pt]
    Accuracy & SQL 查询结果比对得到的执行准确率 \\
    Pre-input Tokens/Q & 每题预输入 token，即可被 token 缓存命中的稳定上下文部分 \\
    Runtime Input Tokens/Q & 每题 runtime 输入 token，即运行中新增且难以缓存的上下文 \\
    Runtime Output Tokens/Q & 每题 runtime 输出 token，即模型返回的工具调用参数、解释文本和最终 SQL \\
    LLM Rounds/Q & 每题串行 LLM 调用轮次 \\
    Total Tokens/Q & 每题总 token 消耗，包含预输入、runtime 输入和 runtime 输出 token \\
    \bottomrule[1.5pt]
  \end{tabularx}
\end{table}

单题 token 消耗可近似分解为：

\begin{equation}
C_{\mathrm{token}} =
C_{\mathrm{preinput}}+
C_{\mathrm{runtime\_input}}+
C_{\mathrm{runtime\_output}} .
\end{equation}

其中 $C_{\mathrm{preinput}}$ 表示系统提示、工具定义、项目 README 等可被 token 缓存命中的稳定前缀；$C_{\mathrm{runtime\_input}}$ 表示每题运行时产生的问题、检索结果、工具返回和历史消息；$C_{\mathrm{runtime\_output}}$ 表示模型生成的工具调用、分析文本和最终 SQL。
\subsection{对比方法}

完整实验计划比较五个方法：

\begin{table}[htbp]\small
  \centering
  \topcaption{主实验对比方法}
  \label{tab:baselines}
  \begin{tabularx}{\textwidth}{llX}
    \toprule[1.5pt]
    \textbf{方法} & \textbf{目录} & \textbf{定位} \\
    \midrule[1pt]
    Bash Agent & \texttt{Bash-Agent} & 无图谱、无专用预处理的通用工具 agent \\
    Pontis & \texttt{Pontis} & 图谱工作空间与工具智能体 \\
    CHESS & \texttt{CHESS} & 检索、schema selection、候选生成和 unit test pipeline \\
    Alpha-SQL & \texttt{Alpha-SQL} & 基于 MCTS 的 test-time search 方法 \\
    DeepEye-SQL & \texttt{DeepEye-SQL} & value retrieval、schema linking、generation、revision、selection pipeline \\
    \bottomrule[1.5pt]
  \end{tabularx}
\end{table}

为保证公平性，所有方法使用统一的大模型和同一 embedding 模型。

\subsection{Pontis 阶段性结果}

根据目前的阶段性测试，Pontis 在 1533 道题目上正确 988 道，错误 539 道，执行错误或系统错误 6 道，整体执行准确率为 64.45\%。


\subsection{错误分析}

从阶段性日志看，Pontis 的错误并非单纯来自模型生成能力不足，而是集中在以下几类：

（1）BIRD SQL 风格偏好。BIRD 中一些题目对百分比、比例、输出列、排序和 \texttt{COUNT DISTINCT} 有较强 benchmark 风格。模型有时会生成业务上合理但与 gold SQL 口径不同的查询。

（2）局部 schema 理解错误。当表名、列名或业务术语相近时，模型可能选择错误字段。例如同名 \texttt{name}、\texttt{id}、\texttt{type} 字段需要结合表语义和样例值消歧。

（3）value grounding 错误。模型需要把自然语言中的实体值映射到数据库中的实际字符串、枚举或编码。若样例值检索不足，容易产生条件偏差。

（4）JOIN 路径偏差。复杂数据库中正确连接路径可能需要桥接表或中间关系。即使模型识别了相关表，也可能遗漏桥接表或连接到错误列。


这些错误类型对应了 Pontis 后续优化方向：加强来自 BIRD train 集的经验复用约束力、改进局部消歧实体、增强 join 路径检查和 value retrieval 能力。

\subsection{消融实验设计}

为证明 Pontis 的模块贡献，本文设计以下消融实验：

\begin{table}[htbp]\small
  \centering
  \topcaption{Pontis 消融实验设计}
  \label{tab:ablation}
  \begin{tabularx}{\textwidth}{llX}
    \toprule[1.5pt]
    \textbf{设置} & \textbf{去掉的能力} & \textbf{验证目标} \\
    \midrule[1pt]
    full & 无 & 全量 Pontis 的准确度与 token 成本 \\
    no-bird-train & 不复用 BIRD train 集生成的经验 & 直接 zero-shot 测试时的准确度和 token 成本变化 \\
    no-guardrail & 禁用护栏机制 & 护栏对 SQL 可靠性、修正轮次和 token 消耗的影响 \\
    \bottomrule[1.5pt]
  \end{tabularx}
\end{table}

其中，\texttt{full} 表示完整 Pontis 设置；\texttt{no-bird-train} 不复用 BIRD train 集生成的经验，直接进行 zero-shot 测试；\texttt{no-guardrail} 保留工具和图谱上下文，但关闭护栏反馈，用于观察护栏对准确度和调用轮次的影响。

\subsection{指标统计协议}

不同方法的成本不宜只用总 token 粗略比较。Pontis 会把稳定系统提示、工具定义和项目级摘要放入可缓存的预输入上下文，而运行时检索结果、工具返回和中间消息属于每题新增上下文；CHESS、Alpha-SQL 和 DeepEye-SQL 则可能通过多阶段候选生成或搜索产生更多运行时输入和输出 token。因此，本文将 token 拆分为预输入、runtime 输入和 runtime 输出三部分。

实验日志按每道题记录最终 SQL、执行比对结果、各轮 LLM 的输入输出 token、可缓存前缀长度和 LLM 调用轮次。对支持 token cache 的模型，预输入 token 单独列出，不与 runtime 输入 token 混合；总 token 消耗则作为统一的资源口径，用于观察不同系统在准确度之外的推理开销。

在预期趋势上，Bash Agent 的预输入 token 较少，但每题 runtime 输入可能因为重复探索而增长；Pontis 的预输入 token 较多，但其中较大部分可被缓存命中，目标是在准确度和 runtime token 成本之间取得平衡；多候选或搜索式方法可能通过更高的 LLM 调用轮次换取更高准确度。本文因此同时报告准确度、token 拆分和 LLM 轮次，而不只比较单一总量。

\section{总结与展望}

% ==================== 后置部分 ====================
\backmatter

\begin{acknowledgment}


\end{acknowledgment}

\bibliography{sample,references}

\appendixs{Pontis 代码模块说明}

Pontis 的主要代码结构如下：

\begin{verbatim}
Pontis/
├── storage/          # Workspace、Neo4j 图接口和发布触发机制
├── extractor/        # 自动提取模块和流水线执行器
├── explorer/         # agent-driven 的关系发现、消歧和摘要模块
├── tool/             # agent-facing 工具实现
├── agent/            # LLM 工具循环、prompt 组装和 guardrails
├── scripts/BIRD/     # BIRD 抽取和 benchmark 脚本
├── docs/             # 架构、性能和实验分析文档
├── pontis_cli.py     # 命令行入口
└── pontis.yml        # 项目与 Neo4j 图配置
\end{verbatim}

本文涉及的主要工具和模块可通过以下命令查看：

\begin{verbatim}
python -m extractor list
pontis <project_path>
python -m scripts.BIRD.run_bird_benchmark --db card_games --limit 5
\end{verbatim}

\end{document}
